\begin{enumerate}
\item Sử dụng định lí: $\widehat{A} + \widehat{B} + \widehat{C} + \widehat{D} = 360^\circ$.
\item Khi biết $ góc, muốn tính góc thứ tư:
80669\widehat{D} = 360^\circ - (\widehat{A} + \widehat{B} + \widehat{C}).80669
\item Khi các góc cho dưới dạng tỉ lệ hoặc chứa ẩn $ (ví dụ $\widehat{A}:\widehat{B}:\widehat{C}:\widehat{D} = 1:2:3:4$):
\begin{itemize}
\item Biểu diễn các góc theo một biến $ (ví dụ $\widehat{A}=x, \widehat{B}=2x, \widehat{C}=3x, \widehat{D}=4x$).
\item Lập phương trình tổng bằng ^\circ$, giải tìm $ rồi suy ra số đo từng góc.
\end{itemize}
\item Sử dụng tính chất góc kề bù để tính góc ngoài hoặc góc trong tương ứng: $\widehat{A} + \widehat{A}_{\text{ngoài}} = 180^\circ$.
\end{enumerate}
